\documentclass[UTF8,a4paper,11pt]{ctexart}

\usepackage[left=2.2cm,right=2.2cm,top=2.1cm,bottom=2.2cm]{geometry}
\usepackage{amsmath,amssymb,mathtools,bm}
\usepackage{booktabs,tabularx,array,multirow}
\usepackage{graphicx}
\usepackage{xcolor}
\usepackage{enumitem}
\usepackage{caption}
\usepackage{hyperref}
\usepackage{fancyhdr}
\usepackage{microtype}
\usepackage{float}

\definecolor{TitleBlue}{RGB}{22,67,105}
\definecolor{RuleBlue}{RGB}{55,105,145}
\definecolor{SoftBlue}{RGB}{238,246,251}
\definecolor{SoftGray}{RGB}{246,247,249}
\definecolor{WarnRed}{RGB}{154,45,45}

\hypersetup{
  colorlinks=true,
  linkcolor=TitleBlue,
  citecolor=TitleBlue,
  urlcolor=TitleBlue,
  pdfauthor={HUAWEI EDA GP project},
  pdftitle={精确非光滑全局布局中的优化器比较与非局部搜索方向研究}
}

\captionsetup{font=small,labelfont=bf,labelsep=quad}
\setlist{nosep,leftmargin=2em}
\setlength{\parindent}{2em}
\setlength{\parskip}{0.25em}
\renewcommand{\arraystretch}{1.18}
\setcounter{secnumdepth}{3}
\setcounter{tocdepth}{2}

\pagestyle{fancy}
\fancyhf{}
\fancyhead[L]{\small 精确非光滑全局布局技术报告}
\fancyhead[R]{\small 2026-09-08}
\fancyfoot[C]{\small\thepage}
\renewcommand{\headrulewidth}{0.4pt}

\DeclareRobustCommand{\code}[1]{\texttt{\detokenize{#1}}}
\newcommand{\R}{\mathbb{R}}
\newcommand{\norm}[1]{\left\lVert #1\right\rVert}
\newcommand{\pos}[1]{\left[#1\right]_{+}}
\newcommand{\rms}{\operatorname{RMS}}
\newcommand{\rmsnorm}{\operatorname{RMSNorm}}
\newcommand{\clip}{\operatorname{clip}}
\newcommand{\fresh}{\textit{fresh exact audit}}

\begin{document}

\begin{center}
  {\LARGE\bfseries\color{TitleBlue} 精确非光滑全局布局中的优化器比较与非局部搜索方向研究\par}
  \vspace{0.45em}
  {\large 基于 \code{adaptec1} 历史中间布局的机制验证与方案设计\par}
  \vspace{0.7em}
  {\normalsize 技术汇报稿（阶段性版本）\quad 2026年9月8日}
\end{center}

\vspace{0.5em}
\noindent\colorbox{SoftBlue}{%
  \parbox{\dimexpr\textwidth-2\fboxsep\relax}{%
  \textbf{报告范围。}本文首先汇报从历史中间结果
  \code{h221_dct_poisson_512/best.pl} 出发，对不同优化器、学习率、坐标步幅裁剪和
  $\lambda$ 调节方式所做的数值比较；随后给出五类 non-local、non-smooth 搜索方向的理论动机、
  数学构造与当前实现边界，并加入五模块在另一份 H221 checkpoint 上的最新 300 步 exact 轨迹。
  该批次用于机制筛选；由于尚缺同骨架 exact-only control，本文只陈述可复核的正负信号，不作超出
  当前 case、checkpoint 和参数的普适性结论。}}

\begin{abstract}
全局布局同时要求降低线长与控制密度溢出，而 exact HPWL 和基于矩形--bin 精确相交的密度量均为
分段光滑甚至局部平坦的函数。本文研究一种保持 exact 评价不变的连续非光滑优化路线：以
epsilon-active HPWL 方向和 exact overlap density 方向构成联合搜索方向，通过动态 $\lambda$ 平衡
线长收缩与密度恢复，并比较 SGD、Heavy Ball、Normalized SGD、AdaGrad、Adam、AMSGrad 与
Dual Averaging 的数值行为。100 步对照表明 AMSGrad 在无逐步接受/回溯的轨迹反馈机制下取得
$99.14\,\mathrm{M}$ HPWL、$5.42\%$ overflow；300 步调参进一步取得
$94.06\,\mathrm{M}$、$6.75\%$，但尚未达到 $86\,\mathrm{M}$ 且 overflow $<7\%$ 的目标。
进一步采用分阶段容量压力调度后，Adam 在 300 步 GP 预算下达到
$87.60\,\mathrm{M}$、$7.00\%$，说明优化器与步幅固然重要，但搜索方向和阶段化任务分解具有更大影响。
基于这一判断，本文提出多尺度 active-set、累计切割压力、盈余--空闲运输、非光滑容量电荷场和
有限半径 exact secant 五类非局部方向，目标是在不引入平滑目标、硬可行域门控或逐步接受机制的条件下，
补足局部次梯度缺失的“远处容量在哪里”与“跨越平台需要移动多远”两类信息。最新 300 步对照显示，
transport 与单探针 finite-radius 分别观察到
$92.21\,\mathrm{M}@6.82\%$ 和 $92.38\,\mathrm{M}@6.60\%$ 的中途可行状态，但终态均重新越过
$7\%$；charge 以 $6.82\%$ 结束，却将 HPWL 恶化至 $120.30\,\mathrm{M}$；cut pressure 未返回
可行带。五种配置均未达到 $86\,\mathrm{M}/<7\%$，说明非局部信息已经产生可用瞬态方向，但长时稳定性、
wire--density 协调和 $\lambda$ 反馈仍是主要瓶颈。
\end{abstract}

\tableofcontents

\section{问题背景与评价契约}

\subsection{研究对象与目标}

设可动单元中心坐标为
\[
  x=(x_1,y_1,\ldots,x_n,y_n)\in\R^{2n}.
\]
本项目研究 ISPD 2005 Bookshelf 数据上的全局布局，不包含 row/site snapping、overlap removal、
legalization 或 detailed placement。因此，“overflow 低于 $7\%$”只表示满足当前全局布局密度指标，
不能表述为布局已经合法化。

本阶段的量化目标为
\[
  W(x)\le 86\times 10^6,
  \qquad O(x)<7\%.
\]
所有实验均从同一个历史中间结果出发，并在 canonical $512\times512$ bins、target density $1.0$
下评价。输入文件经可动单元 die clamp 后的起点为
\[
  W(x_0)=109{,}837{,}622.009,
  \qquad O(x_0)=7.8193\%.
\]
该 checkpoint 的 SHA-256 为
\par\noindent\hspace*{2em}{\small\path{9d0bd99966c3d679780ea73f0c28d9b3eb4163661fe679dfd85d7d32277e79b2}}。

\subsection{Exact HPWL}

对网络 $e$ 中的 pin $p$，pin 坐标包含相对单元中心的偏移：
\[
  X_p=x_{v(p)}+\Delta x_p,
  \qquad
  Y_p=y_{v(p)}+\Delta y_p.
\]
加权 exact HPWL 定义为
\begin{equation}
  W(x)=\sum_{e\in\mathcal E}w_e
  \left[
  \max_{p\in e}X_p-\min_{p\in e}X_p
  +\max_{p\in e}Y_p-\min_{p\in e}Y_p
  \right].
  \label{eq:hpwl}
\end{equation}
式\eqref{eq:hpwl}中的 $\max/\min$ 使目标在 active pin 切换处不可微；多个 pin 并列极值时，
次梯度还具有集合值特征。

\subsection{Exact overlap density 与 overflow}

令 $A_b$ 为 bin $b$ 的面积，$\operatorname{occ}_b(x)$ 为所有可动矩形与该 bin 的精确相交面积，
再加上占据物理容量的固定宏相交面积。定义
\begin{align}
  \rho_b(x)&=\frac{\operatorname{occ}_b(x)}{A_b},\\
  D(x)&=\sum_b\frac{1}{2}A_b\pos{\rho_b(x)-\rho_t}^{2},
  \label{eq:density-energy}\\
  O(x)&=\frac{\sum_b A_b\pos{\rho_b(x)-\rho_t}}
  {A_{\mathrm{movable}}},
  \qquad \pos{z}=\max(z,0).
  \label{eq:overflow}
\end{align}
其中 fixed macro 占用容量但不移动，\code{terminal_NI} 不占物理容量。需要特别强调：
$D(x)$ 是用于构造密度方向的平方超额能量，$O(x)$ 是最终报告的线性超额比例；二者相关但不等价。
因此，$D$ 下降不保证 $O$ 在每一步同步下降。

\subsection{搜索方向与评价指标的分离}

本研究允许使用 epsilon-active、容量粗网格、有限位移探测等信息生成方向，但最终数值均由
式\eqref{eq:hpwl}--\eqref{eq:overflow}重新计算。也就是说，方向 oracle 可以扩展，评价契约不能改变。
历史 DCT/Poisson 阶段仅用于提供起点与对照认识，不进入正式 exact non-smooth 评价链。

\section{联合连续下降框架与 \code{exact_joint_gp}}

\subsection{模块对应的计算环节}

\code{exact_joint_gp} 并不是单一“优化器名称”，而是一段由多个环节组成的联合全局布局阶段：
\begin{enumerate}
  \item \textbf{线长方向：}由 exact HPWL 的 epsilon-active 极值 pin 集合产生
  $g_W\in\partial_{\varepsilon_W}W(x)$；
  \item \textbf{密度方向：}由 exact rectangle/bin overlap 能量产生
  $g_D\in\partial_{\varepsilon_D}D(x)$；
  \item \textbf{动态权衡：}由 $\lambda_k$ 控制线长方向与密度方向的相对强度；
  \item \textbf{节点预条件：}依据 pin 数和单元面积抑制尺度差异；
  \item \textbf{优化器更新：}将联合方向交给 SGD/Adam/AMSGrad 等优化器；
  \item \textbf{边界处理与审计：}移动后仅做 die boundary clamp，并用 exact evaluator 记录轨迹。
\end{enumerate}

当前实现中，第 $i$ 个可动单元的联合方向可概括为
\begin{equation}
  g_{k,i}=
  \frac{g_{W,k,i}+\lambda_k\bigl(g_{D,k,i}+w_Bg_{B,k,i}\bigr)}
  {\max\{1,\,d_i+\lambda_k A_i\}},
  \label{eq:joint-direction}
\end{equation}
其中 $d_i$ 是节点 pin 数，$A_i$ 是单元面积，$g_B$ 表示可选 net-batch 密度共享方向。
V5/V6 单阶段调参关闭逐步 batch acceptance，因此轨迹更新为
\begin{equation}
  x_{k+1}=\Pi_{\mathrm{die}}
  \left(x_k-\Delta_k\right),
  \qquad
  \Delta_k=\mathcal O_k(g_k;\eta,m),
  \label{eq:unconditional-update}
\end{equation}
其中 $\mathcal O_k$ 是具体优化器，$\eta$ 为学习率，$m\eta$ 为单坐标最大更新量。
$\Pi_{\mathrm{die}}$ 只处理物理边界，不是 overflow 可行性投影。

\paragraph{“bridge”的含义。}
\code{exact_joint_gp bridge} 指采用较弱密度压力、允许 overflow 暂时上升的联合 GP 阶段。
它连接“在较宽容量带内快速压缩 HPWL”的探索阶段与“提高密度压力、返回 $7\%$”的 retighten 阶段。
bridge 能显著降低 HPWL，并非因为更换了最终评价函数，而是因为它暂时减少了密度方向对线长收缩的阻碍；
随后必须由 retighten/recovery 重新支付密度代价。其本质是阶段化压力调度，而不是神秘的单步改进算子。

\subsection{动态 \texorpdfstring{$\lambda$}{lambda}：软反馈而非硬约束}

对纯下降实验，可给定从较高探索区向目标区收缩的参考轨迹。令
\[
  p_k=\frac{k}{T-1},\qquad s(p)=p^2(3-2p),
\]
则
\begin{equation}
  \bar O_k=O_{\mathrm{end}}+
  (O_{\mathrm{start}}-O_{\mathrm{end}})[1-s(p_k)],
  \qquad O_{\mathrm{start}}=15\%,\quad O_{\mathrm{end}}=7\%.
  \label{eq:overflow-reference}
\end{equation}
基于误差
\[
  e_k=\frac{O_k-\bar O_k}{O_{\mathrm{scale}}}
\]
可在 log-domain 中更新
\begin{equation}
  \log\lambda_{k+1}=clip
  \left(
  \log\lambda_k+\eta_pe_k+\eta_d(e_k-e_{k-1})+\eta_iI_k
  \right),
  \qquad I_k=\beta_I I_{k-1}+e_k.
  \label{eq:lambda-pid}
\end{equation}
当 overflow 高于参考值时，$\lambda$ 自然增大；低于参考值时，线长方向获得更多自由度。
这里的 $15\%$ 是反馈参考区，不是 hard ceiling；$O_k>\bar O_k$ 不会触发拒绝、回溯或恢复旧布局。

\section{第一部分工作：优化器与步幅机制的数值比较}

\subsection{所比较的优化器}

设 $g_k$ 为式\eqref{eq:joint-direction}的联合方向，以下给出主要更新规则。SGD 与 Heavy Ball 分别为
\begin{align}
  \Delta_k&=\eta g_k,\\
  v_k&=\mu v_{k-1}+(1-\mu)g_k,
  \qquad \Delta_k=\eta v_k.
\end{align}
AdaGrad 采用逐坐标累计平方梯度
\begin{equation}
  r_k=r_{k-1}+g_k\odot g_k,
  \qquad
  \Delta_k=\eta\frac{g_k}{\sqrt{r_k}+\epsilon}.
\end{equation}
Adam 采用一、二阶指数矩估计
\begin{align}
  m_k&=\beta_1m_{k-1}+(1-\beta_1)g_k,\\
  v_k&=\beta_2v_{k-1}+(1-\beta_2)g_k\odot g_k,\\
  \Delta_k&=\eta\frac{\widehat m_k}{\sqrt{\widehat v_k}+\epsilon}.
\end{align}
AMSGrad 进一步令 $\overline v_k=\max(\overline v_{k-1},\widehat v_k)$，避免有效学习率因二阶矩下降而再次增大。
Dual Averaging 则依据累计方向 $G_k=\sum_{t=1}^{k}g_t$ 更新，当前实现使用
$\Delta_k=\eta G_k/\sqrt{k}$。所有优化器最终都应用统一的坐标裁剪：
\begin{equation}
  \Delta_{k,j}\leftarrow\clip(\Delta_{k,j},-m\eta,m\eta).
  \label{eq:step-clipping}
\end{equation}

\subsection{V5：100 步优化器筛选}

V5 固定 \code{exact_joint_gp}、trajectory $\lambda$、$\eta=0.002$、$m=4$、100 步、threads=1，
关闭逐步 acceptance 与 backtracking，仅改变优化器。表\ref{tab:v5}中的 best-feasible 是运行末端
selector 在内存中选取的 $O\le7\%$ 状态；该 selector 用于最终报告/交接，不参与中间步的接受判断。

\begin{table}[!htbp]
\centering
\caption{V5：相同连续下降协议下的优化器比较}
\label{tab:v5}
\small
\begin{tabular}{lrrrr}
\toprule
优化器 & best-feasible HPWL & overflow & 首次 $\le7\%$ & $\lambda$ 峰值 \\
\midrule
AMSGrad & \textbf{99,136,915.82} & 5.4217\% & 54 & 51 \\
Adam & 99,960,810.50 & 5.5886\% & 56 & 51 \\
AdaGrad & 101,850,109.32 & 6.9879\% & 75 & 130 \\
Dual Averaging & 105,949,370.15 & 6.9907\% & 96 & 131 \\
SGD & 108,191,594.03 & 5.2690\% & 10 & 3.2 \\
Heavy Ball & 108,228,173.46 & 5.2480\% & 11 & 3.3 \\
Normalized SGD & 108,985,009.69 & 4.9752\% & 21 & 3.4 \\
\bottomrule
\end{tabular}
\end{table}

结果说明，在当前高度异质、活跃集合频繁切换的梯度上，逐坐标自适应缩放比原始统一步长更有效。
AMSGrad 和 Adam 分别较起点降低约 $9.75\%$ 和 $9.01\%$ HPWL；SGD 类方法虽然很快回到
$7\%$ 以下，却因实际移动较小而没有充分利用 HPWL 下降空间。AdaGrad 与 Dual Averaging 的
$\lambda$ 峰值超过 130，说明后半程大量预算用于密度恢复，而不是继续压缩线长。

\IfFileExists{../framework/results/analysis/20260908_h221traj_optimizer_screen/h221traj_convergence.png}{%
\begin{figure}[!htbp]
  \centering
  \includegraphics[width=0.92\textwidth]{../framework/results/analysis/20260908_h221traj_optimizer_screen/h221traj_convergence.png}
  \caption{V5 七种优化器的 HPWL、overflow 与 $\lambda$ 轨迹。曲线用于比较动态过程，最终数值以 exact audit 表格为准。}
  \label{fig:v5-convergence}
\end{figure}}{}

\subsection{V6：300 步学习率、步幅裁剪与 \texorpdfstring{$\lambda$}{lambda} 调参}

V6 将预算扩展到 300 步，目标直接设为 $86\,\mathrm{M}$ 且 overflow $<7\%$。
主要研究变量为学习率 $\eta$、maximum delta multiplier $m$、$\lambda$ policy 及更新间隔。
表\ref{tab:v6-step}汇总了能够解释主要结论的代表性配置。

\begin{table}[!htbp]
\centering
\caption{V6：代表性学习率与坐标步幅裁剪结果（AMSGrad + ratio $\lambda$）}
\label{tab:v6-step}
\small
\begin{tabular}{lrrrrr}
\toprule
配置 & $\eta$ & $m$ & best-feasible HPWL & overflow & 峰值 overflow \\
\midrule
lr-$5\times10^{-4}$ & 0.0005 & 4 & 97,837,619.00 & 6.9776\% & 13.7142\% \\
lr-$10^{-3}$ & 0.0010 & 4 & 96,491,160.77 & 6.9431\% & 18.8673\% \\
lr-$2\times10^{-3}$ & 0.0020 & 4 & 95,827,835.14 & 6.9621\% & 22.4159\% \\
lr-$4\times10^{-3}$ & 0.0040 & 4 & 99,386,752.93 & 6.9810\% & 23.9030\% \\
\midrule
$m=2$ & 0.0020 & 2 & 95,237,671.77 & 6.9001\% & 22.4159\% \\
$m=1$ & 0.0020 & 1 & 95,226,587.83 & 6.9014\% & 22.4160\% \\
$m=0.5$ & 0.0020 & 0.5 & 94,746,041.05 & 6.9084\% & 19.4579\% \\
$m=0.25$ & 0.0020 & 0.25 & \textbf{94,061,868.14} & \textbf{6.7474\%} & 13.7913\% \\
$m=0.125$ & 0.0020 & 0.125 & 94,188,027.39 & 6.8695\% & 9.7796\% \\
\bottomrule
\end{tabular}
\end{table}

最优配置为 AMSGrad、$\eta=0.002$、$m=0.25$、$(\beta_1,\beta_2)=(0.90,0.99)$，
并使用 ratio $\lambda$、update interval $2$、density weight scale $0.5$。其 exact 结果为
\[
  W=94{,}061{,}868.14,
  \qquad O=6.7474\%.
\]
相对起点 HPWL 下降 $14.36\%$，相对 V5 的 100 步 AMSGrad 结果进一步下降 $5.12\%$；
但距离 $86\,\mathrm{M}$ 仍差 $8.06\,\mathrm{M}$，所以目标尚未命中。

V6 反映出三个重要机制。第一，$\eta$ 不是越大越好，$0.004$ 产生更高早期 excursion 并恶化终值。
第二，在 $\eta=0.002$ 下，将 $m$ 从 4 缩小到 0.25 显著降低早期 overflow 冲击，使更多迭代真正用于
HPWL 下降；但继续缩至 0.125 后收益反转，说明过强裁剪会造成进展不足。第三，trajectory $\lambda$
将 overflow 压到约 $4.7\%$--$4.8\%$，峰值 $\lambda$ 可达 81--412，却没有把额外容量裕量转化为
更低 HPWL；ratio $\lambda$ 保持在目标带附近反而更优。这说明瓶颈不只是密度压力大小，更是方向质量。

\IfFileExists{../framework/results/analysis/20260908_h221_86m_tuning/convergence_main.png}{%
\begin{figure}[!htbp]
  \centering
  \includegraphics[width=0.96\textwidth]{../framework/results/analysis/20260908_h221_86m_tuning/convergence_main.png}
  \caption{V6 代表性配置的收敛曲线。步幅裁剪减少了早期 overflow 过冲，并提高后期 HPWL 改善效率。}
  \label{fig:v6-convergence}
\end{figure}}{}

\subsection{V7：分阶段压力调度对优化器与步长结论的修正}

V6 表明单一 \code{exact_joint_gp} 阶段难以达到 $86\,\mathrm{M}$。V7 因而采用
\[
  \text{cap-15 recovery}
  \rightarrow \text{bridge}(120)
  \rightarrow \text{retighten}(180)
  \rightarrow \text{cap-7 polish}\times2,
\]
其中 bridge 使用较低密度压力，retighten 使用较高密度压力，总 GP 预算仍为 300 步。
表\ref{tab:v7-opt}给出不同优化器在同一链条下的最终 exact 指标。

\begin{table}[!htbp]
\centering
\caption{V7：分阶段链中的优化器比较，bridge/retighten 学习率为 $0.002/0.001$}
\label{tab:v7-opt}
\small
\begin{tabular}{lrrl}
\toprule
优化器 & final HPWL & final overflow & 是否达到 $88\,\mathrm{M}/7\%$ \\
\midrule
Adam & \textbf{87,599,730.55} & 6.9999995\% & 是 \\
AMSGrad & 89,827,947.94 & 6.9999998\% & 否 \\
Heavy Ball & 93,534,597.25 & 10.1388\% & 否 \\
SGD & 93,552,350.68 & 10.1229\% & 否 \\
Normalized SGD & 94,203,472.85 & 9.9623\% & 否 \\
AdaGrad & 94,335,504.44 & 6.9999999\% & 否 \\
Dual Averaging & 94,374,897.03 & 6.9999982\% & 否 \\
\bottomrule
\end{tabular}
\end{table}

Adam 的学习率呈现明显单峰。bridge/retighten 学习率依次取
$0.001/0.0005$、$0.002/0.001$ 与 $0.004/0.002$ 时，分别得到
$89.26\,\mathrm{M}$、$87.60\,\mathrm{M}$ 和 $88.61\,\mathrm{M}$。最佳链在各阶段的 exact audit 为
\[
\begin{aligned}
109.84\,\mathrm{M}@7.82\%
&\rightarrow96.92\,\mathrm{M}@15.00\%\\
&\rightarrow87.10\,\mathrm{M}@15.00\%\\
&\rightarrow90.33\,\mathrm{M}@6.86\%\\
&\rightarrow87.60\,\mathrm{M}@7.00\%.
\end{aligned}
\]
这一轨迹清楚展示了 bridge 的作用：先以容量裕量换取大幅 HPWL 压缩，再由 retighten 与 exact recovery
返回目标密度带。V7 达到 $88\,\mathrm{M}/7\%$，但仍未达到本文最终的 $86\,\mathrm{M}/<7\%$ 目标。

\IfFileExists{../framework/results/analysis/20260908_h252_h257_chain/optimizer_screen.png}{%
\begin{figure}[!htbp]
  \centering
  \includegraphics[width=0.78\textwidth]{../framework/results/analysis/20260908_h252_h257_chain/optimizer_screen.png}
  \caption{V7 分阶段链的优化器终态比较。能否在 retighten 预算内返回 $7\%$ 是区分优化器的重要条件。}
  \label{fig:v7-optimizer}
\end{figure}}{}

终态散点只能回答“最后到达哪里”，不能显示不同优化器在 bridge 与 retighten 切换时的瞬态行为。
图\ref{fig:v7-convergence}进一步给出完整 300 步 exact 轨迹：前 120 步允许以较高 overflow 换取
HPWL 压缩，第 121 步切入 retighten 后，各方法先出现一次压力重置引起的 excursion，随后向 $7\%$
目标带恢复。Adam 的优势并非全程单调下降，而是能够在较低 HPWL 区域完成后期密度恢复。

\IfFileExists{../framework/results/analysis/20260908_h252_h257_chain/convergence_hpwl.png}{%
\begin{figure}[!htbp]
  \centering
  \includegraphics[width=0.90\textwidth]{../framework/results/analysis/20260908_h252_h257_chain/convergence_hpwl.png}
  \vspace{0.5em}
  \includegraphics[width=0.90\textwidth]{../framework/results/analysis/20260908_h252_h257_chain/convergence_overflow.png}
  \caption{V7 优化器与学习率实验的完整收敛曲线：exact HPWL（上）与 exact overflow（下）。横轴按全局 GP 步编号；第 1--120 步为 bridge，第 121--300 步为 retighten，圆点表示阶段边界的 fresh exact audit。}
  \label{fig:v7-convergence}
\end{figure}}{}

\subsection{阶段性判断}

上述结果支持如下判断：
\begin{enumerate}
  \item 自适应逐坐标缩放能够显著改善 exact active-set 方向的可用性，但最优优化器依赖阶段结构；
  单阶段纯下降中 AMSGrad 最优，分阶段 bridge/retighten 链中 Adam 最优。
  \item 学习率与坐标裁剪需要联合解释。$\eta$ 决定全局尺度，$m$ 控制局部极端坐标；适度裁剪可以避免
  早期 active-set 突变消耗大量密度恢复预算。
  \item 仅调高 $\lambda$ 不能解决方向缺少目的地信息的问题。过大的 $\lambda$ 会得到低 overflow，
  却可能浪费可用于降低 HPWL 的容量裕量。
  \item 从 $94.06\,\mathrm{M}$ 到 $87.60\,\mathrm{M}$ 的主要收益来自阶段化机制，而非局部超参数微调。
  因此后续重点应从“继续大规模扫 optimizer”转向“构造更有信息的非局部非光滑方向”。
\end{enumerate}

\section{第二部分工作：五类 Non-local Non-smooth 方向方案}

\subsection{共同动机：局部次梯度为什么容易失效}

Exact overlap density 对单元位置是分段低阶函数。当单元在当前 overlap 几何区域内移动时，局部梯度可能
很小甚至为零；只有跨过 bin edge 后，密度能量才发生明显变化。于是可能出现
\[
  D(x+\delta)=D(x)\quad (\norm{\delta}\le r_0),
  \qquad
  D(x+\Delta)<D(x)\quad (\norm{\Delta}>r_0).
\]
此外，局部密度排斥通常只表达“离开拥塞 bin”，却不表达“应前往哪个空闲区域”。这与平滑
DCT/Poisson 方法形成鲜明对比：后者通过全局势场把远处空闲容量传播到当前位置，方向稠密且随位置连续变化，
因此早期下降很快；代价是优化目标已经成为平滑 surrogate，且场的改善不等同于 exact overflow 的改善。

五个方案遵循同一原则：\textbf{保留 exact objective 和 exact evaluator，只扩充 direction oracle。}
共同更新骨架为
\begin{equation}
  g_D^{(m)}=\alpha_E\widehat g_D^{\mathrm{exact}}
  +\alpha_A\widehat g_D^{(m,\mathrm{aux})},
  \qquad
  g_k=\rmsnorm\!\left(\widehat g_W+\lambda_k\widehat g_D^{(m)}\right),
  \label{eq:nonlocal-common}
\end{equation}
其中 $m$ 表示具体 oracle。每一步均按式\eqref{eq:unconditional-update}无条件应用；exact audit 仅用于
测量、$\lambda$ 反馈和日志，不触发 hard constraint、candidate acceptance、backtracking、rollback 或
best-feasible restore。五个模块相互独立，不共享 occupancy cache、optimizer moment 或隐式状态。

\subsection{方案一：多尺度 Density \texorpdfstring{$\varepsilon$}{epsilon}-Active Oracle}

\subsubsection{动机}

单一 $\varepsilon=0$ 的 exact density 次梯度只看当前 active geometry。若同时观察邻域内多个潜在 active face，
就可能在真正跨越几何断点之前获得非零方向，从而缓解 zero-subgradient 和过度局部的问题。

\subsubsection{数学构造}

定义 enlarged subdifferential
\begin{equation}
  \partial_{\varepsilon}D(x)=
  \operatorname{conv}
  \left(
  \bigcup_{\norm{z-x}_{\infty}\le\varepsilon}\partial D(z)
  \right).
  \label{eq:epsilon-subdiff}
\end{equation}
实际计算只取有限尺度
\[
  \varepsilon_s\in\{0.25h,0.5h,h,2h\},
\]
其中 $h$ 为 bin 特征尺寸。每个尺度先独立归一化，再进行凸组合：
\begin{equation}
  g_A^{\mathrm{MS}}=
  \frac{1}{|\mathcal S_+|}
  \sum_{s\in\mathcal S_+}
  \frac{g_D^{(\varepsilon_s)}}
  {\rms(g_D^{(\varepsilon_s)})+\epsilon}.
  \label{eq:multiscale}
\end{equation}
由于 $\varepsilon=0$ 的 exact density direction 已在公共项中单独存在，辅助尺度不重复包含它。

\subsubsection{预期作用与边界}

该方法仍由有限个 piecewise exact active-face 方向构成，不定义新的平滑密度目标。其优势是实现轻量、
与 exact overlap 几何关系紧密；局限是信息仍然局部，无法显式指出远处空闲容量的位置。

\subsection{方案二：累计切割压力（Cumulative Cut Pressure）}

\subsubsection{动机}

若某条竖直切割线左侧的单元总面积超过左侧可用容量，则无论局部如何重排，必有一部分面积跨越该切割线。
这是一条来自面积守恒的全局必要条件，可为局部密度方向提供长程信号。

\subsubsection{数学构造}

对竖直 cut $c$，定义左侧面积与容量
\begin{align}
  M_L(c)&=\sum_{b:x_b<c}\operatorname{occ}_b,\\
  C_L(c)&=\sum_{b:x_b<c}\rho_tA_b,
\end{align}
以及累计盈余
\begin{equation}
  S_x(c)=M_L(c)-C_L(c).
  \label{eq:cut-surplus}
\end{equation}
若 $S_x(c)>0$，左侧至少有 $S_x(c)$ 的面积需要向右跨越 cut。可构造分段线性势
\begin{equation}
  H_c(x_i)=\pos{S_x(c)}\pos{c-x_i},
  \label{eq:cut-potential}
\end{equation}
其负次梯度把 cut 左侧单元推向右侧。对 $S_x(c)<0$、水平 cut 以及相反方向作对称处理，叠加所有 cut
即可获得多尺度全局容量压力。

\subsubsection{预期作用与边界}

该方法只需 prefix scan 与 cell lookup，复杂度接近 $O(B+N)$；不使用 FFT、Poisson、Gaussian 或
平滑密度图。其局限在于方向主要沿坐标轴，对障碍物和不规则二维空白区域的表达能力弱于二维运输或电荷场。

\subsection{方案三：盈余--空闲运输方向（Surplus--Deficit Transport）}

\subsubsection{动机}

局部排斥只回答“从哪里离开”，运输模型进一步回答“应该去哪里”。令
\begin{equation}
  s_b=\pos{\operatorname{occ}_b-C_b},
  \qquad
  d_b=\pos{C_b-\operatorname{occ}_b},
  \qquad C_b=\rho_tA_b.
  \label{eq:supply-demand}
\end{equation}
其中 $s_b$ 是源端盈余，$d_b$ 是目标端空闲容量。

\subsubsection{数学构造}

理想离散运输计划为
\begin{equation}
\begin{aligned}
  \min_{\pi\ge0}\quad &\sum_{b,c}\pi_{bc}
  \left(|x_b-x_c|+|y_b-y_c|\right),\\
  \text{s.t.}\quad &\sum_c\pi_{bc}\le s_b,
  \qquad
  \sum_b\pi_{bc}\le d_c.
\end{aligned}
\label{eq:transport}
\end{equation}
源 bin $b$ 的期望位移向量为
\begin{equation}
  v_b=\frac{\sum_c\pi_{bc}(p_c-p_b)}
  {\sum_c\pi_{bc}+\epsilon},
  \qquad g_b^{T}=-v_b.
  \label{eq:transport-direction}
\end{equation}
于是梯度下降 $x\leftarrow x-\eta g^T$ 会将源区单元推向空闲区。

\subsubsection{当前原型边界}

当前轻量原型采用 coarse grid 与最近 deficit 的稀疏匹配思路，并非已经求解式\eqref{eq:transport}的完整
min-cost flow。正式报告结果时应称为“粗粒度盈余--空闲运输方向原型”，并同步记录是否满足质量守恒、
运输质量和平均距离；只有实现完整约束后，才能称为严格最优运输 oracle。

\subsection{方案四：非光滑容量电荷场（Capacity Charge Field）}

\subsubsection{动机}

希望获得类似电势布局的远程 capacity awareness，但不引入 Poisson 方程与光滑电势目标。
将一轮 fresh occupancy 下的容量不平衡冻结为
\begin{equation}
  q_b=\operatorname{occ}_b-C_b.
  \label{eq:charge}
\end{equation}
$q_b>0$ 表示拥塞源，$q_b<0$ 表示空闲容量。

\subsubsection{数学构造}

采用 Manhattan 距离与有限支撑 hinge kernel
\begin{equation}
  r_{ib}=|x_i-c_b^x|+|y_i-c_b^y|,
  \qquad K_R(r)=\pos{R-r}.
\end{equation}
辅助势为
\begin{equation}
  H_C(x\mid q)=\sum_iA_i\sum_bq_bK_R(r_{ib}).
  \label{eq:charge-potential}
\end{equation}
对冻结的 $q$ 取次梯度，正电荷产生排斥，负电荷产生吸引；$r_{ib}\ge R$ 时相互作用为零。
该方法只包含绝对值和 hinge，因而是分段线性的非光滑场，而不是 $1/r$ 电势或谱卷积。

\subsubsection{预期作用与边界}

二维有限支撑场比 cut pressure 更适合表达不规则空闲区，并可通过
$R\in\{2h,4h,8h\}$ 控制作用尺度。主要风险是人工场与 exact overflow 的相关性不足，或大量稀疏方向被
全局 RMS 归一化过度放大。因此需要记录 active RMS、非零比例、符号契约和有效相互作用对数。

\subsection{方案五：有限半径 Exact Coordinate Oracle}

\subsubsection{动机}

当 infinitesimal 次梯度为零时，可直接用有限位移跨越平台。对单元 $i$ 的四个坐标方向
$u\in\{+e_x,-e_x,+e_y,-e_y\}$，探测 exact density energy 的单边 secant：
\begin{equation}
  s_{i,u}^{D}(\delta)=
  \frac{D(x+\delta u)-D(x)}{\delta}.
  \label{eq:finite-secant}
\end{equation}
选择能够降低 $D$ 的方向作为辅助 gradient-like vector，但不直接提交 probe candidate。
实际布局仍由统一优化器按式\eqref{eq:unconditional-update}连续更新。

\subsubsection{断点与计算预算}

优先选择接近下一几何事件的 $\delta$，例如单元边界跨越 bin boundary；对未来的有限半径 HPWL oracle，
则可考虑 pin 进入/离开 net 极值集合的断点。由于 exact probe 需要在短生命周期布局副本上重新评价，
该方案计算量最大，应优先探测局部 wire/exact-density 梯度范数较小的单元：
\begin{equation}
  r_i=\sqrt{g_{W,i,x}^2+g_{W,i,y}^2+g_{D,i,x}^2+g_{D,i,y}^2}.
\end{equation}
修订后的 oracle 支持按低信息程度选点和 $\pm x,\pm y$ 四方向，并令辅助项只使用 density secant，
避免在公共框架外再次混合 HPWL 与 density。由于 exact probe 代价较高，本次 300 步正式记录采用
\code{probe_count=1}、$R=h$ 的快速配置；它验证的是单探针机制，不代表原计划 32 个探针的完整效果。

\subsection{五种方向的互补关系}

\begin{table}[!htbp]
\centering
\caption{五类非局部非光滑方向的作用尺度与核心信息}
\label{tab:five-oracles}
\small
\begin{tabularx}{\textwidth}{>{\raggedright\arraybackslash}p{2.9cm}
>{\raggedright\arraybackslash}p{2.5cm}
>{\raggedright\arraybackslash}X
>{\raggedright\arraybackslash}p{3.2cm}}
\toprule
方案 & 信息尺度 & 核心补充信息 & 主要风险/限制 \\
\midrule
多尺度 active-set & 局部到邻域 & 邻近 active face 与即将出现的几何 kink & 仍不指向远处空闲区 \\
累计切割压力 & 全局一维投影 & 面积守恒下必须跨 cut 的净流量 & 轴对齐，二维障碍表达有限 \\
盈余--空闲运输 & 全局二维 & 从拥塞源到空闲 sink 的目的地 & 匹配质量、守恒和 cell 映射 \\
容量电荷场 & 有限半径二维 & 正盈余排斥、负容量吸引 & 场与 exact overflow 的一致性 \\
有限半径 secant & 单元级有限位移 & 跨越局部平台后的 exact 变化 & 评价开销高，需稀疏选点 \\
\bottomrule
\end{tabularx}
\end{table}

这五类方案并非同一种机制的不同参数：多尺度 active-set 扩大局部广义梯度；cut pressure 编码全局守恒；
transport 显式寻找目的地；charge field 提供有限支撑二维远程作用；finite-radius secant 则跨越
infinitesimal oracle 看不见的平台。独立模块对照的价值正在于分离这些信息来源。

\section{五方案的统一实验与最新数值结果}

\subsection{输入、共同参数与证据边界}

最新批次使用的输入为外部 \code{h221_a1_fine/best.pl}，其 SHA-256 为
\par\noindent\hspace*{2em}{\small\path{3129818214a86ba0ee88d281448d782bfb6285186dd7ad4efafb95bc1342e348}}。
在 canonical $512\times512$、target density $1.0$ 和 movable clamp 后，统一起点为
\[
  W(x_0)=109{,}425{,}511.1896,
  \qquad O(x_0)=7.803537\%.
\]
该 SHA 与 V5/V6 使用的历史 \code{h221_dct_poisson_512/best.pl} 不同，故两组结果只能作数值参照，
不能作为严格的单变量因果比较。

五个模块均运行 300 步，固定 AMSGrad、$\eta=0.002$、maximum delta $m=0.25$、
$(\beta_1,\beta_2)=(0.90,0.99)$、exact/auxiliary density weights $1.0/0.1$，并采用
$15\%\rightarrow7\%$ trajectory $\lambda$、update interval $2$、$\lambda\in[0.01,256]$。
evaluator、threads=1 和 metrics-only retention 完全一致。每个更新均无条件应用，只有 die clamp；
日志中的 accepted=300/rejected=0 表示 300 个 applied steps，并不表示存在候选接受门控。

模块特有参数分别为：Multi-scale 使用
$\{0.25h,0.5h,h,2h\}$，Transport 使用 $32\times32$ coarse bins，Charge 使用 $R=4h$，
Finite-radius 使用 $R=h$ 且 \code{probe_count=1}。运行元数据记录的源码状态为
\code{e230425 + dirty}；对应的方向符号与尺度修复随后提交为 \code{a5e3042}，当前结果图排除了未完成的
多探针 finite 运行。当前批次尚未包含同骨架的 exact-only control，因此表\ref{tab:nonlocal-main}只能回答
“五个当前配置各自表现如何”，不能单独量化 auxiliary oracle 相对于 $\alpha_A=0$ 的净贡献。

\subsection{终态与轨迹最优值}

Pure descent 的 stage 输出是最后一个 applied state，不会恢复中途最优布局。因此表\ref{tab:nonlocal-main}
同时列出 best-any、best-feasible-observed 与 last。best-any 若对应 $O>7\%$，只能用于解释探索能力；
best-feasible-observed 是离线轨迹统计，不是算法内部 acceptance 或 restore。

\begin{table}[!htbp]
\centering
\caption{五个 non-local 模块的 300 步 exact 轨迹摘要；HPWL 单位为 $10^6$}
\label{tab:nonlocal-main}
\scriptsize
\begin{tabular}{lrrrrrrr}
\toprule
方案 & best-any & 对应 ovf. & best-feas. & 对应 ovf. & 迭代 & last HPWL & last ovf. \\
\midrule
Multi-scale active & 89.842 & 8.6038\% & 95.744 & 6.8274\% & 128 & 126.923 & 8.9225\% \\
Cut pressure & 91.876 & 26.6971\% & -- & -- & -- & 183.712 & 13.1667\% \\
Transport & 91.943 & 25.7410\% & \textbf{92.211} & 6.8214\% & 255 & 95.436 & 8.8449\% \\
Charge field & \textbf{88.777} & 10.7042\% & 113.346 & 6.9976\% & 282 & 120.303 & \textbf{6.8218\%} \\
Finite-radius ($p=1$) & 91.873 & 25.9219\% & 92.383 & \textbf{6.5954\%} & 266 & \textbf{93.504} & 7.3593\% \\
\bottomrule
\end{tabular}
\end{table}

\begin{table}[!htbp]
\centering
\caption{五个模块的恢复时刻、峰值 excursion 与终态审计}
\label{tab:nonlocal-audit}
\scriptsize
\begin{tabular}{lrrrrrr}
\toprule
方案 & 首次 $\le7\%$ & peak ovf. & final $D$ & final $\rho_{\max}$ & runtime/s & final HPWL 变化 \\
\midrule
Multi-scale active & 102 & 32.7698\% & 485,454.92 & 2.3777 & 73.85 & $+15.99\%$ \\
Cut pressure & -- & 35.3097\% & 628,188.18 & 2.1042 & 30.91 & $+67.89\%$ \\
Transport & 120 & 35.1855\% & 381,614.48 & 2.3364 & 31.50 & $-12.78\%$ \\
Charge field & 282 & 31.9383\% & 353,987.37 & 2.2752 & 43.83 & $+9.94\%$ \\
Finite-radius ($p=1$) & 123 & 35.1993\% & 234,742.34 & 2.0255 & 101.30 & $-14.55\%$ \\
\bottomrule
\end{tabular}
\end{table}

五种轨迹在前 17--19 步均出现 $31.94\%$--$35.31\%$ 的 overflow 峰值。这并非 hard constraint
失效，而是当前 pure-descent 协议允许的真实 excursion；但峰值远高于 $15\%$ soft reference，说明
$\lambda$ 反馈和第一阶段步幅仍偏慢。随后各方法呈现明显分化。

\IfFileExists{../framework/results/analysis/20260908_h221_nonlocal_fixed_final/hpwl.png}{%
\begin{figure}[!htbp]
  \centering
  \includegraphics[width=0.80\textwidth]{../framework/results/analysis/20260908_h221_nonlocal_fixed_final/hpwl.png}
  \vspace{0.4em}
  \includegraphics[width=0.80\textwidth]{../framework/results/analysis/20260908_h221_nonlocal_fixed_final/overflow.png}
  \caption{五个模块的 exact HPWL（上）与 exact overflow（下）轨迹。虚线为 $7\%$，点线为
  $15\%$ soft reference；所有曲线均为无接受、无回滚的 300 步连续更新。}
  \label{fig:nonlocal-convergence}
\end{figure}}{}

\subsection{分模块结果解释}

\paragraph{Multi-scale active。}
该方向在第 102 步首次回到 $7\%$，并于第 128 步观察到
$95.744\,\mathrm{M}@6.827\%$，相对本批次起点降低 $12.50\%$ HPWL；但第 193 步的最低 HPWL
$89.842\,\mathrm{M}$ 已不可行，随后轨迹反向恶化，最终为 $126.923\,\mathrm{M}@8.923\%$。
这说明多尺度 active face 含有瞬态有效信息，但当前固定权重与长时动量不能维持稳定交换。

\paragraph{Cut pressure。}
Cut 在 300 步内从未回到 $7\%$，终态 HPWL 增至 $183.712\,\mathrm{M}$，是本批次明确的负结果。
其最小 HPWL 出现在高达 $26.70\%$ overflow 的第 39 步；随后累计切割压力与 wire 方向长期失配。
该结果否定的是当前轴向聚合、权重和控制配置，尚不能外推为所有 cut-based capacity oracle 均无效。

\paragraph{Transport。}
Transport 在第 120 步首次回到 $7\%$，并在第 255 步观察到全批次最低可行 HPWL
$92.211\,\mathrm{M}@6.821\%$，相对起点降低 $15.73\%$。其终态仍回升到
$95.436\,\mathrm{M}@8.845\%$，说明“源到 sink”的目的地信息能够产生有效可行点，但当前 controller
无法把恢复后的轨迹锁定在目标带内。

\paragraph{Charge field。}
Charge 给出全批次最低 best-any HPWL $88.777\,\mathrm{M}$，但对应 overflow 为 $10.704\%$。
它到第 282 步才首次可行，best-feasible-observed 为 $113.346\,\mathrm{M}$，终态虽达到
$6.822\%$，HPWL 却比起点高 $9.94\%$。该现象表明容量吸引/排斥能完成密度恢复，但与 wire 方向的
后期协调不足，恢复代价过高。

\paragraph{Finite-radius。}
单探针配置在第 266 步观察到 $92.383\,\mathrm{M}@6.595\%$，终态
$93.504\,\mathrm{M}@7.359\%$，具有全批次最低 last HPWL，但仍未满足最终 overflow 条件。
其运行时间为 101.30 s，约为 Transport 的 3.2 倍；即使每轮只探测一个低信息单元，exact secant
成本仍然显著。该结果提供了跨越局部平台的正信号，同时说明扩大 probe count 前必须先优化评价预算。

\IfFileExists{../framework/results/analysis/20260908_h221_nonlocal_fixed_final/hpwl_overflow_phase.png}{%
\begin{figure}[!htbp]
  \centering
  \includegraphics[width=0.76\textwidth]{../framework/results/analysis/20260908_h221_nonlocal_fixed_final/hpwl_overflow_phase.png}
  \caption{五个模块在 $(O,W)$ 平面中的相轨迹。Transport 与 finite-radius 多次穿越 $7\%$ 邻域但未稳定停留；
  Charge 完成密度回收却沿高 HPWL 分支结束；Cut 则持续进入高代价区域。}
  \label{fig:nonlocal-phase}
\end{figure}}{}

\subsection{本批次能够与不能支持的结论}

当前结果支持两点。其一，non-local oracle 并非完全没有下降信息：Transport 与 finite-radius 均产生了
约 $92.2\,\mathrm{M}$ 的中途可行状态，明显优于各自起点。其二，主要困难已经从“是否能找到下降方向”
转变为“能否在无接受机制下稳定留在 $7\%$ 附近”；最终只有 Charge 可行，但其 HPWL 更差。

当前结果不能支持以下结论：
\begin{enumerate}
  \item 不能把 $88.777\,\mathrm{M}@10.704\%$ 称为正式最优结果，也不能把离线统计的中途可行点当作模块输出；
  \item 不能因 Cut 的单一配置失败而否定全部切割方法；
  \item 不能把单探针 Finite-radius 结果外推到 32 探针；
  \item 不能在缺少同骨架 exact-only control 时宣称收益全部来自 auxiliary oracle。
\end{enumerate}
下一轮应优先补充 $\alpha_A=0$ 对照、早期步幅/$\lambda$ 响应消融，并在 telemetry 中持久化
$\lambda_k$、step max、clipped fraction、auxiliary active RMS 与 wire--density cosine，以解释当前的
大幅循环轨迹。

\section{讨论：为何平滑化方法下降快，而局部非光滑方法容易停滞}

\subsection{信息传播范围不同}

局部 exact density 次梯度由当前重叠的 bin 边界决定，其感受野有限。远处存在大块空闲区时，局部方向通常
无法感知。Poisson/DCT 方法求解全局势场，相当于将全芯片容量不平衡传播到每个位置，因而几乎所有可动单元
都获得稠密方向。这是平滑方法早期下降快的首要原因，而不是 exact 目标在数学上“不可优化”。

\subsection{平台、active-set 抖动与动量失配}

在 piecewise 区域内部，exact 次梯度可能为零；跨越边界后又会突变。统一步长因此面临两难：步长小则长期
停在平台，步长大则同时跨越大量 active face，造成 overflow 冲击。动量还可能在 active-set 切换后继续沿旧方向
推进。V6 中适度 maximum-delta 裁剪有效，正是因为它限制了少数极端坐标，而没有令所有坐标同步缩步。

\subsection{目标冲突与标量 \texorpdfstring{$\lambda$}{lambda} 的局限}

HPWL 方向常将单元拉向网络几何中心，密度方向则将其推出拥塞区。单一标量 $\lambda$ 只控制总体强弱，
不能解决不同区域、不同单元上的方向冲突，也不能创造不存在的目的地信息。V6 中 trajectory $\lambda$ 将
overflow 压得很低却没有达到 86M，说明“更强密度惩罚”不等于“更好的 HPWL--overflow 交换”。

\subsection{平滑化方法的优势与边界}

平滑化带来连续、稠密、低噪声的方向，适合大步推进与 Nesterov/Adam 类加速；但其 surrogate 的下降并不保证
exact HPWL 或 exact overflow 同步下降，并且平滑尺度可能掩盖细粒度几何断点。因此本项目的研究方向不是简单
排斥平滑信息，而是尝试用 non-local、non-smooth、可解释 oracle 获得类似的长程信息，同时保持 exact evaluator
作为唯一结论依据。

\section{结论与下一步工作}

本文阶段性工作可归纳为四点。第一，无逐步 acceptance 的连续非光滑下降机制能够从 H221 起点返回 $7\%$
密度带；AMSGrad 在 V5 的 100 步单阶段筛选中最佳。第二，V6 的 300 步调参表明学习率 $0.002$ 与
maximum delta $0.25$ 的组合显著优于默认大步幅，但单阶段最优值 $94.06\,\mathrm{M}$ 仍未达到 86M。
第三，V7 的 bridge/retighten/recovery 分阶段链将结果推进到 $87.60\,\mathrm{M}@7.00\%$，说明阶段结构
比继续扩大超参数扫描更重要。第四，五个 non-local 模块的最新 300 步结果进一步表明，Transport 和
单探针 Finite-radius 能产生 $92.2\,\mathrm{M}$ 量级的中途可行状态，但没有以可行终态结束；Charge 能返回
$7\%$ 内却牺牲 HPWL，Cut 在当前配置下失败。五种配置均未达到 $86\,\mathrm{M}/<7\%$。

因此，下一阶段不宜直接扩大五模块融合或进行大规模 optimizer 扫描。优先事项是补齐同一 common engine 下的
exact-only control，校准最初 20 步的 step/$\lambda$ 响应，并记录 auxiliary 方向稀疏度、尺度与
wire--density cosine。若 Transport/Finite-radius 在该对照下仍保持正信号，再研究自适应 auxiliary weight、
更快的 finite probe 和不依赖硬接受门控的阻尼机制。所有后续结论仍应同时报告 last、best-feasible-observed
和 best-any，避免将不可行轨迹点或离线最优点冒充 stage 输出。

\begin{thebibliography}{9}
\bibitem{adam}
D. P. Kingma and J. Ba,
``Adam: A Method for Stochastic Optimization,''
\textit{International Conference on Learning Representations}, 2015.

\bibitem{amsgrad}
S. J. Reddi, S. Kale, and S. Kumar,
``On the Convergence of Adam and Beyond,''
\textit{International Conference on Learning Representations}, 2018.

\bibitem{dualavg}
Y. Nesterov,
``Primal-dual subgradient methods for convex problems,''
\textit{Mathematical Programming}, vol. 120, pp. 221--259, 2009.

\bibitem{eplace}
J. Lu, H. Zhuang, P. Chen, H. Chang, C.-C. Chang, Y.-C. Wong, L. Sha, D. Huang, Y. Luo,
C.-C. Teng, and C.-K. Cheng,
``ePlace: Electrostatics-based placement using fast Fourier transform and Nesterov's method,''
\textit{Proceedings of the 51st ACM/EDAC/IEEE Design Automation Conference}, 2014.
\end{thebibliography}

\end{document}
